% GENERATED FILE. Edit canonical lessons, then run npm run generate:books.
% canonical-source-hash: fnv1a64:479c803fc2eb4c0a
% canonical-lessons: RU-C343-kritikovat, RU-C343-kritika, RU-C343-odobryat, RU-C343-osuzhdat, RU-C343-uprekat

\chapter{Criticising and approving}
\label{ch:ru-criticising-and-approving}
\hlchaptermodality{343}

\begin{chapteropening}
\textbf{By the end of this chapter:} \emph{I can criticise and take criticism, say that I approve or disapprove of something, and reproach someone.}

Complete the last lesson of chapter 343: I can criticise and take criticism, say that I approve or disapprove of something, and reproach someone.
\end{chapteropening}

\section[\texorpdfstring{\ru{критиковать} (kritikovát)}{критиковать (kritikovát)}]{\ru{критиковать} (kritikovát) — to criticise}

\label{lesson:RU-C343-kritikovat}



\begin{quote}
Before the new one: say the Russian for to speak one’s mind, then the Russian for to share.
\end{quote}

\subsection*{You'll want to know: \ru{критиковать}}
\textbf{\ru{критиковать}} (\emph{kritikovát}) — \textquotedblleft{}to criticise\textquotedblright{}.

\begin{quote}
\ru{Мне} \ru{не} \ru{нравится}, \ru{когда} \ru{меня} \ru{критикуют} \ru{при} \ru{всех}.
\end{quote}

\begin{quote}
\emph{Mne ne nrávitsya, kogdá menyá kritikúyut pri vsekh.}
\end{quote}

\begin{quote}
I do not like being criticised in front of everyone.
\end{quote}

Imperfective, with the present \emph{\ru{критикую}, \ru{критикуешь}}; perfective \emph{\ru{раскритиковать}}, to tear apart. It takes the accusative.

\begin{cousinweb}
\textbf{\ru{критиковать}} is built on \emph{\ru{критика}}, which is usually traced, through western languages, to Greek \emph{kritik\'{\={e}}} \textquotedblleft{}the art of judging\textquotedblright{}, from \emph{krínein} \textquotedblleft{}to separate, to judge\textquotedblright{}.
\end{cousinweb}

\begin{grammarlens}[title={What you've built}]
One more word for telling, explaining and describing at B1.
\end{grammarlens}

\subsection*{Your turn}
\begin{itemize}
\raggedright
  \item \emph{Say it:} \emph{kritikovát}
  \item \emph{Say it:} \emph{kritikovát}, once more
  \item \emph{Say it:} a sentence of your own with \emph{kritikovát}
  \item \emph{Recall:} say the Russian for to speak one’s mind, then the Russian for to share, then say \emph{kritikovát} again
\end{itemize}

\begin{tcolorbox}[breakable,colback=teal!4,colframe=teal!35!black,title={Before you move on}]
What does \ru{критиковать} mean? (\textquotedblleft{}To criticise\textquotedblright{}.) Say it once more.
\end{tcolorbox}
\section[\texorpdfstring{\ru{критика} (krítika)}{критика (krítika)}]{\ru{критика} (krítika) — criticism}

\label{lesson:RU-C343-kritika}



\begin{quote}
Before the new one: say the Russian for to share, then the Russian for to criticise.
\end{quote}

\subsection*{You'll want to know: \ru{критика}}
\textbf{\ru{критика}} (\emph{krítika}) — \textquotedblleft{}criticism\textquotedblright{}.

\begin{quote}
\ru{Он} \ru{спокойно} \ru{слушает} \ru{критику}.
\end{quote}

\begin{quote}
\emph{On spokóyno slúshayet krítiku.}
\end{quote}

\begin{quote}
He listens calmly to criticism.
\end{quote}

Feminine. \emph{\ru{Конструктивная} \ru{критика}} — constructive criticism. It is also reviewing as a field: \emph{\ru{литературная} \ru{критика}}.

\begin{cousinweb}
\textbf{\ru{критика}} is usually traced, through French and German, to Greek \emph{kritik\'{\={e}}} \textquotedblleft{}the art of judging\textquotedblright{}, built on \emph{krínein} \textquotedblleft{}to separate, to decide\textquotedblright{}.
\end{cousinweb}

\begin{grammarlens}[title={What you've built}]
One more word for telling, explaining and describing at B1.
\end{grammarlens}

\subsection*{Your turn}
\begin{itemize}
\raggedright
  \item \emph{Say it:} \emph{krítika}
  \item \emph{Say it:} \emph{krítika}, once more
  \item \emph{Say it:} a sentence of your own with \emph{krítika}
  \item \emph{Recall:} say the Russian for to share, then the Russian for to criticise, then say \emph{krítika} again
\end{itemize}

\begin{tcolorbox}[breakable,colback=teal!4,colframe=teal!35!black,title={Before you move on}]
What does \ru{критика} mean? (\textquotedblleft{}Criticism\textquotedblright{}.) Say it once more.
\end{tcolorbox}
\section[\texorpdfstring{\ru{одобрять} (odobryát)}{одобрять (odobryát)}]{\ru{одобрять} (odobryát) — to approve of}

\label{lesson:RU-C343-odobryat}



\begin{quote}
Before the new one: say the Russian for to criticise, then the Russian for criticism.
\end{quote}

\subsection*{You'll want to know: \ru{одобрять}}
\textbf{\ru{одобрять}} (\emph{odobryát}) — \textquotedblleft{}to approve of\textquotedblright{}.

\begin{quote}
\ru{Отец} \ru{не} \ru{одобряет} \ru{мой} \ru{выбор}.
\end{quote}

\begin{quote}
\emph{Otéts ne odobryáyet moy výbor.}
\end{quote}

\begin{quote}
My father does not approve of my choice.
\end{quote}

Imperfective; perfective \emph{\ru{одобрить}}. It takes the accusative.

\begin{cousinweb}
\textbf{\ru{одобрять}} is built from \emph{\ru{о}-} and \emph{\ru{добрый}} \textquotedblleft{}good\textquotedblright{}, so to approve is to call a thing good. \emph{\ru{Добрый}} is inherited from Proto-Slavic \emph{*dobrŭ} (EtymDB).
\end{cousinweb}

\begin{grammarlens}[title={What you've built}]
One more word for telling, explaining and describing at B1.
\end{grammarlens}

\subsection*{Your turn}
\begin{itemize}
\raggedright
  \item \emph{Say it:} \emph{odobryát}
  \item \emph{Say it:} \emph{odobryát}, once more
  \item \emph{Say it:} a sentence of your own with \emph{odobryát}
  \item \emph{Recall:} say the Russian for to criticise, then the Russian for criticism, then say \emph{odobryát} again
\end{itemize}

\begin{tcolorbox}[breakable,colback=teal!4,colframe=teal!35!black,title={Before you move on}]
What does \ru{одобрять} mean? (\textquotedblleft{}To approve of\textquotedblright{}.) Say it once more.
\end{tcolorbox}
\section[\texorpdfstring{\ru{осуждать} (osuzhdát)}{осуждать (osuzhdát)}]{\ru{осуждать} (osuzhdát) — to condemn, to disapprove of}

\label{lesson:RU-C343-osuzhdat}



\begin{quote}
Before the new one: say the Russian for criticism, then the Russian for to approve of.
\end{quote}

\subsection*{You'll want to know: \ru{осуждать}}
\textbf{\ru{осуждать}} (\emph{osuzhdát}) — \textquotedblleft{}to condemn, to disapprove of\textquotedblright{}.

\begin{quote}
\ru{Я} \ru{не} \ru{осуждаю} \ru{его}: \ru{он} \ru{очень} \ru{устал}.
\end{quote}

\begin{quote}
\emph{Ya ne osuzhdáyu yegó: on óchen ustál.}
\end{quote}

\begin{quote}
I do not judge him harshly: he is very tired.
\end{quote}

Imperfective; perfective \emph{\ru{осудить}}. In court it also means to convict.

\begin{cousinweb}
\textbf{\ru{осуждать}} is built from \emph{\ru{о}-} and \emph{\ru{судить}} \textquotedblleft{}to judge\textquotedblright{}, on \emph{\ru{суд}}, inherited from Proto-Slavic \emph{*sǫdŭ} (EtymDB). The \emph{\ru{жд}} is the Church Slavonic softening of \emph{\ru{д}}.
\end{cousinweb}

\begin{grammarlens}[title={What you've built}]
One more word for telling, explaining and describing at B1.
\end{grammarlens}

\subsection*{Your turn}
\begin{itemize}
\raggedright
  \item \emph{Say it:} \emph{osuzhdát}
  \item \emph{Say it:} \emph{osuzhdát}, once more
  \item \emph{Say it:} a sentence of your own with \emph{osuzhdát}
  \item \emph{Recall:} say the Russian for criticism, then the Russian for to approve of, then say \emph{osuzhdát} again
\end{itemize}

\begin{tcolorbox}[breakable,colback=teal!4,colframe=teal!35!black,title={Before you move on}]
What does \ru{осуждать} mean? (\textquotedblleft{}To condemn, to disapprove of\textquotedblright{}.) Say it once more.
\end{tcolorbox}
\section[\texorpdfstring{\ru{упрекать} (uprekát)}{упрекать (uprekát)}]{\ru{упрекать} (uprekát) — to reproach}

\label{lesson:RU-C343-uprekat}



\begin{quote}
Before the new one: say the Russian for to approve of, then the Russian for to condemn.
\end{quote}

\subsection*{You'll want to know: \ru{упрекать}}
\textbf{\ru{упрекать}} (\emph{uprekát}) — \textquotedblleft{}to reproach\textquotedblright{}.

\begin{quote}
\ru{Жена} \ru{упрекает} \ru{мужа} \ru{в} \ru{том}, \ru{что} \ru{он} \ru{мало} \ru{бывает} \ru{дома}.
\end{quote}

\begin{quote}
\emph{Zhená uprekáyet múzha v tom, chto on málo byváyet dóma.}
\end{quote}

\begin{quote}
The wife reproaches her husband for being at home so little.
\end{quote}

Imperfective; perfective \emph{\ru{упрекнуть}}. \emph{\ru{Упрекать} \ru{кого} \ru{в} \ru{чём}}. The noun is \emph{\ru{упрёк}}, a reproach.

\begin{cousinweb}
\textbf{\ru{упрекать}} is built on the noun \emph{\ru{упрёк}} \textquotedblleft{}a reproach\textquotedblright{}. The origin of the root is not settled.
\end{cousinweb}

\begin{grammarlens}[title={What you've built}]
One more word for telling, explaining and describing at B1.
\end{grammarlens}

\subsection*{Your turn}
\begin{itemize}
\raggedright
  \item \emph{Say it:} \emph{uprekát}
  \item \emph{Say it:} \emph{uprekát}, once more
  \item \emph{Say it:} a sentence of your own with \emph{uprekát}
  \item \emph{Recall:} say the Russian for to approve of, then the Russian for to condemn, then say \emph{uprekát} again
\end{itemize}

\begin{tcolorbox}[breakable,colback=teal!4,colframe=teal!35!black,title={Before you move on}]
What does \ru{упрекать} mean? (\textquotedblleft{}To reproach\textquotedblright{}.) Say it once more.
\end{tcolorbox}
